\documentclass[11pt, a4paper]{article}

% ── Packages ──────────────────────────────────────────────────────────
\usepackage[utf8]{inputenc}
\usepackage[T1]{fontenc}
\usepackage{lmodern}
\usepackage[margin=1in]{geometry}
\usepackage{amsmath, amssymb, amsthm}
\usepackage{booktabs}
\usepackage{longtable}
\usepackage{graphicx}
\usepackage{hyperref}
\usepackage{xcolor}
\usepackage{listings}
\usepackage{tcolorbox}
\usepackage{polya}

\hypersetup{
  colorlinks=true,
  linkcolor=polyablue,
  urlcolor=polyablue,
  citecolor=polyablue,
}

% ── Code listing style ────────────────────────────────────────────────
\lstdefinestyle{polya-python}{
  language=Python,
  basicstyle=\ttfamily\small,
  keywordstyle=\color{polyablue}\bfseries,
  commentstyle=\color{polyagray}\itshape,
  stringstyle=\color{polyagreen},
  numbers=left,
  numberstyle=\tiny\color{polyagray},
  numbersep=8pt,
  frame=single,
  framerule=0.4pt,
  rulecolor=\color{polyagray},
  breaklines=true,
  showstringspaces=false,
  tabsize=4,
}
\lstset{style=polya-python}
\title{Inspector Assignment}
\author{uber-polya}
\date{2026-02-26}

\begin{document}

\maketitle
\tableofcontents
\newpage

% ── Phase A: Problem Formulation ─────────────────────────────────────
\section{Problem Formulation}

\begin{polyamodel}

\textbf{Problem Type}: Find \hfill
\textbf{Domain}: Graph Theory / Combinatorics

\textbf{Named Problem}: Assignment Problem

\end{polyamodel}

% ── Universe ──────────────────────────────────────────────────────────
\subsection{Universe}

\begin{itemize}
  \item $Assignments: \{Alice, Bob, Carol, Dave, Eve, Frank\}$
\end{itemize}

% ── Variables ─────────────────────────────────────────────────────────
\subsection{Variables}

\begin{longtable}{lll}
\toprule
\textbf{Symbol} & \textbf{Meaning} & \textbf{Type / Range} \\
\midrule
\endhead
$x_{ij}$ & assignment decision & $\{0, 1\}$ \\
$n$ & number of binary variables = 60 & $\mathbb{{Z}}^+$ \\
\bottomrule
\end{longtable}

% ── Structure ─────────────────────────────────────────────────────────
\subsection{Structure}

Assign 6 food safety inspectors to 10 facilities (5 types x 2 each) to
maximize total expertise score. Each inspector handles at most 3 facilities;
each facility needs exactly 1 inspector.

% ── Mapping ───────────────────────────────────────────────────────────
\subsection{Real-World Mapping}

\begin{longtable}{ll}
\toprule
\textbf{Real-World Concept} & \textbf{Mathematical Object} \\
\midrule
\endhead
assignments & $x: solution vector (assignments)$ \\
total score/cost & $objective function value$ \\
\bottomrule
\end{longtable}

% ── Constraints ───────────────────────────────────────────────────────
\subsection{Constraints}

\begin{enumerate}
  \item $\sum_j x_{ij} \le 3$
  \hfill \textit{(capacity limit)}
  \item $\sum_j x_{ij} = 1$
  \hfill \textit{(exact assignment)}
  \item $\sum_i x_{ij} = 1$
  \hfill \textit{(coverage requirement)}
  \item $x_{ij} \in \{0, 1\} \quad \forall\, i, j$
  \hfill \textit{(binary decision variables)}
\end{enumerate}

% ── Objective / Claim ─────────────────────────────────────────────────
\subsection{Objective}

\begin{equation}
\max \sum_{i,j} c_{ij} \, x_{ij}
\end{equation}

% ── Approach ──────────────────────────────────────────────────────────
\subsection{Approach}

\begin{description}
  \item[Suggested approach:] ILP (PuLP/CBC, Branch \& Bound)
  \item[Complexity class:] 
  \item[Available tools:] ILP
\end{description}
% ── Phase B: Solution ────────────────────────────────────────────────
\section{Solution}

\begin{polyaresult}

\textbf{Answer}: 6-element assignments: Alice -> [F1, F2], Bob -> [F3, F4], Carol -> [F5, F6], Dave -> [F7, F8], Eve -> [F9, F10], Frank -> []

\textbf{Objective Value}: $90.0$

\textbf{Optimal}: Yes \hfill
\textbf{Feasible}: Yes
\textbf{Algorithm}: ILP (PuLP/CBC, Branch \& Bound) \quad $$

\textbf{Time}: 0.007788196962792426s

\textbf{Certificate}: LP relaxation = ILP solution (gap = 0\%). Global optimum proven.

\end{polyaresult}

% ── Solution Details ──────────────────────────────────────────────────
\subsection{Solution Details}

Assignments:
  Alice -> F1, F2
  Bob -> F3, F4
  Carol -> F5, F6
  Dave -> F7, F8
  Eve -> F9, F10
  Frank -> 

Objective value: z* = 90.0
LP relaxation: 90.0000
Integrality gap: 0.00\%
Solution efficiency: 100.0\%

Method: Integer Linear Programming via PuLP/CBC.

- 60 binary variables (6 inspectors x 10 facilities)
- 16 constraints (10 facility-coverage + 6 inspector-capacity)
- The LP relaxation is integral (totally unimodular constraint matrix),
  so LP optimal = ILP optimal.

% ── Verification ──────────────────────────────────────────────────────
\subsection{Verification}

\begin{longtable}{lcl}
\toprule
\textbf{Check} & \textbf{Status} & \textbf{Detail} \\
\midrule
\endhead
Feasibility & \passmark & All constraints satisfied \\
Optimality & \passmark & LP relaxation = 90.00, gap = 0.00\% \\
Efficiency & \passmark & 100.0\% of theoretical maximum \\
\bottomrule
\end{longtable}

% ── Phase C: Interpretation ──────────────────────────────────────────
\section{Interpretation}

% ── The Question & The Answer ─────────────────────────────────────────
\subsection{The Question}
Assign 6 food safety inspectors to 10 facilities (5 types x 2 each) to
maximize total expertise score.

\subsection{The Answer}

\begin{polyainsight}[title={Bottom Line}]
The optimal solution has objective value z* = 90.0.
\end{polyainsight}

% ── What This Means ───────────────────────────────────────────────────
\subsection{What This Means}

- Optimal score: 90/90 (100\% efficiency)
- Each facility assigned to its top expert
- 0\% optimality gap

Integer Linear Programming via PuLP/CBC.

- 60 binary variables (6 inspectors x 10 facilities)
- 16 constraints (10 facility-coverage + 6 inspector-capacity)
- The LP relaxation is integral (totally unimodular constraint matrix),
  so LP optimal = ILP optimal.

% ── Sensitivity Analysis ──────────────────────────────────────────────
\subsection{Sensitivity Analysis}

\begin{longtable}{llrrl}
\toprule
\textbf{Parameter} & \textbf{Current} & \textbf{Change} & \textbf{New Obj.} & \textbf{Class} \\
\midrule
\endhead
Alice & present & removed & 88.0 &
\textcolor{polyagray}{sensitive} \\
Bob & present & removed & 85.0 &
\textcolor{polyared}{critical} \\
Carol & present & removed & 85.0 &
\textcolor{polyared}{critical} \\
Dave & present & removed & 86.0 &
\textcolor{polyagray}{sensitive} \\
Eve & present & removed & 87.0 &
\textcolor{polyagray}{sensitive} \\
Frank & present & removed & 90.0 &
\textcolor{polyagreen}{robust} \\
\bottomrule
\end{longtable}

% ── Figures ───────────────────────────────────────────────────────────

% ── Recommendations ───────────────────────────────────────────────────
\subsection{Recommendations}

\begin{enumerate}
  \item The solution is provably optimal; no further improvement is possible within this model.
  \item Monitor critical parameters: Bob, Carol.
\end{enumerate}

% ── Limitations ───────────────────────────────────────────────────────
\subsection{Limitations}

\begin{polyawarnbox}
\begin{itemize}
  \item Model assumes deterministic parameters; real-world variability not captured.
  \item Solution quality depends on the accuracy of input data.
\end{itemize}
\end{polyawarnbox}

% ── Appendix ─────────────────────────────────────────────────────────

\end{document}